Potential outline for introduction:
\begin{itemize}
    \item Define \ac{hmlv} manufacturing (what it is, why it matters).
    \item Highlight challenges: Variability, short product life cycles, small batches.
    \item State survey purpose: e.g., "to review the technological landscape of \ac{hmlv} manufacturing across industry 3.0, 4.0, and 5.0 paradigms."
    \item Mention scope: Not a proposal, but a descriptive and comparative mapping of technologies.
\end{itemize}

Think of the \ac{hmlv} manufacturing field as a stack comprised of three layers.
The foundation, the digital layer, and the human-centric future. With this division, the foundation is what we in this survey present as "Industry 3.0".
It can also be called the foundational techniques which has shaped the beginning of \ac{hmlv} manufacturing.
As the digital side made advances, technological innovations were applied to the field of manufacturing, thereby creating the digital layer comprised of Industry 4.0 technologies such as Cyber-Physical Systems (CPS), Internet of Things (\ac{iot}), Big Data, Artificial Intelligence (AI), and Automation.
The future might lie in the more human-centric, as a counter-reaction to the non-human paradigm of Industry 4.0.
Therefore, the last stack "Human-Centric Future" deals with the future of \ac{hmlv} manufacturing in the case of Industry 5.0 values, such as a more human-centric, resilient and sustainable industrial paradigm.

In this survey, we aim to collect and present the three different paradigms, and what the cutting-edge techniques within each paradigm are.
